% RESULTADO_RECUPERADO: integral dimensional functor, eq:cZ-internos,
% eq:longitud-ruta, and definition of ell_0; composition with Article III.
% The consolidated proof and its algebraic check are specific to this contribution.
\subsection{Constitutive speed and kinematic realization of the clock}
\label{v:sec:enlace-velocidad-constitutiva}

Radiative propagation uses the speed section produced by the
angular transport of the same HMT state. The identification is established
before representation units are chosen. Let $x=A_{\mathrm{rad}}$ and
$y=C^*_{\mathrm{rad}}$ be the angular coordinates generated in
Section~\ref{sec:ii-angulos}, with the channels
\eqref{eq:ii-canales}, and let
$\mathsf R$ be the self-adjoint involution of the two sheets, with rank-one
projectors $P_+$ and $P_-$. On the domain $x>|y|$, the transport and its
channels are
\begin{equation}
 T=\exp(-xI-y\mathsf R)=q_+P_++q_-P_-,
 \qquad q_+=e^{-x-y},\quad q_-=e^{-x+y},\quad 0<q_\pm<1.
 \label{v:eq:velocidad-transporte}
\end{equation}
The response of the temporal blocks of lengths $90$ and $120$ is
\begin{equation}
 \mathsf V=(I-T^{90})(I-T^{120})^{-1}
       =r_+P_++r_-P_-,\qquad
 r_\pm=\frac{1-q_\pm^{90}}{1-q_\pm^{120}}.
 \label{v:eq:velocidad-respuesta}
\end{equation}
The construction starts from the inherited channels; the symbols $x,y$ and
$q_\pm$ denote their coordinates, not parameters fitted to a target
speed.

\begin{proposition}[Identity of the constitutive and kinematic realizations]
\label{v:prop:identidad-velocidad}
In the same dimensional chart, the clock section $c_{\mathrm{int}}$
and the constitutive-response section $c_{\mathrm{vac}}$ coincide:
\begin{equation}
 c_{\mathrm{int}}=c_{\mathrm{vac}}
   =\widehat c\,u_C,
 \qquad
 \widehat c=(r_+r_-)^{-1}
            =\exp(-\mathcal E),
 \label{v:eq:identidad-velocidad}
\end{equation}
where
\begin{equation}
 \mathcal E=
 \log\!\bigl[(1-q_+^{90})(1-q_-^{90})\bigr]
 -\log\!\bigl[(1-q_+^{120})(1-q_-^{120})\bigr].
 \label{v:eq:velocidad-lector-simetrico}
\end{equation}
The equality preserves sheet exchange.
\end{proposition}
\begin{proof}
The factors appearing in the logarithms are positive. Therefore,
\[
 \mathcal E=\log(r_+r_-)
            =\log\det\mathsf V,
 \qquad \exp(-\mathcal E)=(\det\mathsf V)^{-1}.
\]
The determinant is computed here in the two-sheet realization with rank-one
projectors. The kinematic definition $c_{\mathrm{int}}=\exp(-\mathcal E)u_C$
and the constitutive definition $c_{\mathrm{vac}}=(r_+r_-)^{-1}u_C$
therefore give the same section. Exchanging $P_+\leftrightarrow
P_-$ interchanges $r_+$ and $r_-$ and preserves their product. Finally,
$\det T=e^{-2x}$ corresponds to the elementary transport, whereas
$\det\mathsf V=r_+r_-$ corresponds to its temporal response; the two
determinants have distinct roles in this composition.
\end{proof}

\subsubsection{Trajectory length and elementary duration}

The positive dimensional frame $(u_S,u_C,u_T,u_Q)$ induces the length
basis $u_L=u_Cu_T$. For an enriched trajectory $\gamma$ of
$n=|\gamma|$ transitions with a constant channel, the additive readers are
\begin{equation}
 T_{\mathrm{int}}(\gamma)=n u_T,
 \qquad L_{\mathrm{int}}(\gamma)=n\widehat c\,u_L.
 \label{v:eq:velocidad-longitud-reloj}
\end{equation}
Their additivity concerns the count and realization of each transition:
enriched routes additionally preserve orientation and memory.

\begin{corollary}[Normalization of the elementary length]
\label{v:cor:longitud-elemental-constitutiva}
For $n>0$, $L_{\mathrm{int}}(\gamma)/T_{\mathrm{int}}(\gamma)
=c_{\mathrm{vac}}$. In the normalization $t_0=u_T$,
\begin{equation}
 \ell_0=c_{\mathrm{int}}t_0=\widehat c\,u_L.
 \label{v:eq:longitud-elemental-constitutiva}
\end{equation}
In a general temporal chart $t_0=\tau_0u_T$, with $\tau_0>0$,
the corresponding expression is $\ell_0=\tau_0\widehat c\,u_L$.
\end{corollary}
\begin{proof}
Dividing the expressions in \eqref{v:eq:velocidad-longitud-reloj} cancels the integer
$n$ and uses $u_L/u_T=u_C$. The elementary formula is obtained
by multiplying the speed section by $t_0$.
\end{proof}

Generalization to an edge-dependent channel preserves the additive
formula $L_{\mathrm{int}}(\gamma)=\sum_{e\in\gamma}\widehat c(e)u_L$.
Its quotient by $nu_T$ is the mean of the speeds of those edges.
The homogeneous radiative realization uses the constant channel specified
in \eqref{v:eq:velocidad-longitud-reloj}.

\subsubsection{Change of chart and adapted normalization}

\begin{proposition}[Covariance of the identification]
\label{v:prop:covariancia-identidad-velocidad}
Let $u_C'=d u_C$ and $u_T'=\eta u_T$, with $d,\eta>0$.
The same speed and the same elementary duration have coordinates
\begin{equation}
 \widehat c'=d^{-1}\widehat c,
 \qquad \tau_0'=\eta^{-1}\tau_0,
 \qquad u_L'=d\eta u_L.
 \label{v:eq:velocidad-cambio-carta}
\end{equation}
In particular, if $c_V=\widehat c_Vu_C^V$ and
$u_C^V=d u_C^{III}$, equality of the sections is exactly equivalent to
$d\widehat c_V=\widehat c_{III}$.
\end{proposition}
\begin{proof}
Express each section in the two bases. This gives
$\widehat c'u_C'=\widehat c u_C$ and
$\tau_0'u_T'=\tau_0u_T$. Their product satisfies
\[
 \tau_0'\widehat c' u_L'
 =\frac{\tau_0}{\eta}\frac{\widehat c}{d}
       d\eta u_L
 =\tau_0\widehat c u_L=\ell_0.
\]
The final equivalence follows by substituting the relation between the bases.
\end{proof}

The adapted choice $d=\widehat c$ gives $u_C'=c_{\mathrm{vac}}$ and
$\widehat c'=1$. This choice expresses the already constructed section in a
basis adapted to it. The spectral coefficient in the internal chart remains
$(r_+r_-)^{-1}$; spectral inversion uses that internal chart.
For example, retaining the action and charge bases, the constitutive
coordinates transform as
\[
 \widehat\mu'=d\widehat\mu,
 \qquad \widehat\varepsilon'=d\widehat\varepsilon,
 \qquad
 \widehat\mu'\widehat\varepsilon'=(\widehat c')^{-2}.
\]
In the adapted chart, $\widehat\mu'=r_-/r_+$ and
$\widehat\varepsilon'=r_+/r_-$, since the internal coordinates are
$\widehat\mu=r_-^2$ and $\widehat\varepsilon=r_+^2$.
